\section{Hypocrisy and Anglo-Saxon Politics}

\begin{quotex}
It may be remarked that the harsh opinion of Machiavelli has been more widespread in England and the United States than in the nations of Continental Europe. This is no doubt natural, because the distinguishing quality of Anglo-Saxon politics has always been hypocrisy, and hypocrisy must always be at pains to shy away from the truth. It is also the case that judgments of Machiavelli are usually based upon acquaintance with \emph{The Prince} alone, an essay which, though plain enough, can be honestly misinterpreted when read out of the context of the rest of his writings. However, something more fundamental than these minor difficulties is at stake.

\flright{\textsc{James Burnham}, \textit{The Machiavellians}\footnote{\url{https://books.google.com/books?id=rGU6AAAAMAAJ\&q=machiavellians\&dq=machiavellians\&pgis=1}}}
\end{quotex}

\flrightit{Posted on 2008-08-21 by Cologero}

\begin{center}* * *\end{center}

\begin{footnotesize}
\begin{sffamily}
\texttt{Attila on 2013-06-02 at 12:27 said:}

I trust Anglo-Saxons (professionally and personally) as far as I can throw them.

\hfill

\texttt{Helmut Schillinger on 2016-12-20 at 01:39 said:}

Add to this that despite the blunt evil the the German nazis represented, the hypocricy of the Brits again showed itself in the brutal war crimes of their beloved hero Churchill, who against the council of his generals, decided to wipe out a city populated only by the old, children and women. Genocide in India, China and Africa in their very special racist and back handed way is also the trademark of Anglosaxon hypocricy. To this day the Brits still act as if they were the great Empire that once cut its bloody trail throughout world history. They still adore their medieval aristocracy even though their economy is in shambles.

\hfill

\texttt{J on 2017-08-21 at 16:10 said:}

Mr. Cologero, I don't mean to be rude, but were you jilted by an Anglo-Saxon woman or something? This is one of at least a half dozen ``knocks'' on Anglo-Saxons on this site.

\hfill

\texttt{Cologero on 2017-08-21 at 21:16 said:}

No, J, but I've jilted several Anglo-Saxon women.

A ``knock'' is gratuitous, without foundation. So there are no ``knocks'' on this site.

Perhaps you should address this: \textit{Short Note on Woman in East and West}\footnote{\url{http://www.gornahoor.net/?p=6420}}.

\hfill

\texttt{J on 2017-08-22 at 00:44 said:}

``In regard to the woman, the Islamic system was complete, whole. In Europe, the situation is horrendous: man pushes woman into the street, removes her veil, prostitutes her in sunlight, places her where she cannot be, where it is against nature for her to be: in the schools, trains, cafes, in the pigsty. The Anglo-Saxon peoples wanted that: and among us, with the stupidity that pushes us toward the ordure, we are following the example of those coarse beasts.''

If everyone else is stupid enough to follow along, the Anglo-Saxons should be praised for their leading role!

I just found it interesting, is all. With all the Evola-citing alone, there are plenty of ethnic groups you could -- since you don't like ``knock'' -- ``critique.'' Is it because those are without foundation? Or, like Batman, can the Anglo-Saxons ``take it''?

\hfill

\end{sffamily}
\end{footnotesize}

